\section{Kondensatorlarni parallel v aketma ket ulash}


\begin{enumerate}

\item 100 nF, 200 nF va 500 nF sig‘imli uchta kondensatorni ketma-ket ulangandagi umumiy sig‘imini toping (nF).
\begin{enumerate}[label=\Alph*)]
\item 800
\item 500
\item 100
\item 58,8
\end{enumerate}

\item Elektr sig‘imlari bir xil bo‘lgan uchta kondensator 2 V kuchlanish manbaiga parallel ulandi. Bunda har bir kondensator 1 nC zaryadga ega bo‘ldi. Zanjirning umumiy sig‘imini toping (pF).
\begin{enumerate}[label=\Alph*)]
\item 1400
\item 1500
\item 1600
\item 800
\end{enumerate}

\item Elektr sig‘imi 6 $\mu F$, 2 $\mu F$ va 4 $\mu F$ ga teng bo‘lgan kondensatorlar 24 V kuchlanish manbaiga ketma-ket ulangan. Ikkinchi kondensator qoplamalari orasidagi potensiallar farqini toping (V).
\begin{enumerate}[label=\Alph*)]
\item 8
\item 12
\item 4
\item 13,1
\end{enumerate}

\item O‘zgarmas tok manbaiga ulangan 8 $\mu F$ sig‘imli kondensatorga parallel ulangan voltmetr 6 V kuchlanishni ko‘rsatdi. Agar bu kondensatorga 2 $\mu F$ sig‘imli ikkinchi kondensator ketma-ket ulansa, voltmetr qanday kuchlanishni (V) ko‘rsatadi?
\begin{enumerate}[label=\Alph*)]
\item 1,2
\item 2
\item 6
\item 4,6
\end{enumerate}

\item $C_1 = 1 \, \mu F$ va $C_2 = 3 \, \mu F$ sig‘imli ikki kondensatorni ketma-ket ulab hosil qilingan elektr sxema $U = 220 \, V$ li o‘zgarmas kuchlanish manbaiga ulangan. Har bir kondensatordagi kuchlanishni aniqlang (V).
\begin{enumerate}[label=\Alph*)]
\item 115; 105
\item 120; 100
\item 145; 65
\item 165; 55
\end{enumerate}

\item 3 $\mu F$ va 6 $\mu F$ sig‘imli kondensatorlar o‘zaro ulangan va 120 V EYUK li batareyaga biriktirilgan. Agar kondensatorlar ketma-ket ulangan bo‘lsa, ularning qoplamalari orasidagi potensiallar farqi aniqlansin (V).
\begin{enumerate}[label=\Alph*)]
\item 80; 40
\item 80; 80
\item 40; 40
\item 80; 20
\end{enumerate}

\item  Rasmda tasvirlangan kondensatorlar tizimining umumiy sig‘imini toping.

\begin{circuitikz}[scale=0.8, transform shape]
  % Kirish terminali va 4C kondensator
  \draw (0,0) to [short, o-] (0.5,0)
        to [C, l=$4C$] (2,0);
  
  % Birinchi tugun va 3C kondensatorlar (Parallel)
  \draw (2,0) -- (2,1)
        to [C, l=$3C$] (4,1);
  \draw (2,0) -- (2,-1)
        to [C, l=$3C$] (4,-1);
        
  % 3C larning chiqishini bitta nuqtaga birlashtirish
  \draw (4,1) -- (4,-1);
  
  % Ikkinchi tugun va uchta 2C kondensatorlar (Parallel)
  \draw (4,1) to [C, l=$2C$] (6,1);
  \draw (4,0) to [C, l=$2C$] (6,0);
  \draw (4,-1) to [C, l=$2C$] (6,-1);
  
  % Barcha 2C larning chiqishini birlashtirish va chiqish terminali
  \draw (6,1) -- (6,-1);
  \draw (6,0) to [short, -o] (7,0);

  % Nuqtalarni belgilash (ixtiyoriy)
  \node[circle, fill, inner sep=1pt] at (2,0) {};
  \node[circle, fill, inner sep=1pt] at (4,0) {};
  \node[circle, fill, inner sep=1pt] at (6,0) {};
\end{circuitikz}

\begin{enumerate}[label=\Alph*)]
\item $47C/16$
\item $14C/15$
\item $15C/7$
\item $12C/7$
\end{enumerate}

\item  Sxemadan umumiy sig‘imni toping.
\begin{enumerate}[label=\Alph*)]
\item $3C/2$
\item $3C$
\item $6C$
\item $2C$
\end{enumerate}

\item  Rasmda har bir kondensatorning sig‘imi C ga teng. M va N nuqtalar orasidagi potensiallar farqi 120 V bo‘lsa, K va L nuqtalar orasidagi potensiallar farqini (V) toping.
\begin{enumerate}[label=\Alph*)]
\item 60
\item 40
\item 20
\item 30
\end{enumerate}

\item  Sig‘imlari $C_1 = 3 \, \mu F, C_2 = 5 \, \mu F, C_3 = 6 \, \mu F$ va $C_4 = 5 \, \mu F$ bo‘lgan to‘rtta kondensator rasmda ko‘rsatilgan sxema bo‘yicha ulangan. Batareyaning umumiy sig‘imini aniqlang ($\mu F$).
\begin{enumerate}[label=\Alph*)]
\item 2
\item 3,4
\item 9
\item 4,5
\end{enumerate}

\item  Rasmda ko‘rsatilgan kondensatorlar batareyasining A va B nuqtalar orasidagi elektr sig‘imini aniqlang.
\begin{enumerate}[label=\Alph*)]
\item 0,4C
\item 0,2C
\item C
\item 0,8C
\end{enumerate}

\item  Rasmdagi kondensatorlar batareyasining sig‘imi 5,8 $\mu F$ ga teng. Agar $C_2 = 1 \, \mu F, C_3 = 4 \, \mu F$ bo‘lsa, birinchi kondensatorning sig‘imi ($\mu F$) va undagi zaryad qanday ($\mu C$)?
\begin{enumerate}[label=\Alph*)]
\item 5; 1,1
\item 10; 1,1
\item 5; 0,8
\item 1,4; 3
\end{enumerate}

\item  Rasmdagi zanjirda K va L uchlar orasidagi zaryad miqdori Q ga teng. MN uchlari orasiga ulangan C sig‘imli kondensatorning zaryadi nimaga teng?
\begin{enumerate}[label=\Alph*)]
\item $Q/5$
\item $Q/2$
\item $8Q/53$
\item $13Q/40$
\end{enumerate}

\item Rasmda tasvirlangan kondensatorlar tizimining umumiy sig‘imini toping.
\begin{enumerate}[label=\Alph*)]
\item 3C
\item 6C/11
\item 5C/11
\item 6C
\end{enumerate}

\item  Ketma-ket ulangan ikkita bir xil kondensator modda bilan ikki xil usulda yarmigacha to‘ldirilgan. Bunda kondensator kuchlanishlarining farqi qanday bo‘ladi (V)?
\begin{enumerate}[label=\Alph*)]
\item 0
\item 24
\item 28,5
\item 12,6
\end{enumerate}

\item  Rasmda ko‘rsatilgan kondensatorlar batareyasida a va b nuqtalar orasidagi kuchlanishini toping (V).
\begin{enumerate}[label=\Alph*)]
\item 0
\item 8,25
\item 7,75
\item 5,33
\end{enumerate}

\item Sig‘imi 3 $\mu F$ zaryadi 5 $\mu C$ va sig‘imi 8 $\mu F$ zaryadi 3 $\mu C$ bo‘lgan kondensatorlar mos qutblari bilan ulangandan keyin, birinchi kondensatordagi zaryadni toping ($\mu C$).
\begin{enumerate}[label=\Alph*)]
\item 24/11
\item 64/11
\item 6/11
\item 16/11
\end{enumerate}

\item Sig‘imi 3 $\mu F$ zaryadi 5 $\mu C$ va sig‘imi 8 $\mu F$ zaryadi 3 $\mu C$ bo‘lgan kondensatorlar mos qutblari bilan ulangandan keyin, ikkinchi kondensatordagi zaryadni toping ($\mu C$).
\begin{enumerate}[label=\Alph*)]
\item 24/11
\item 64/11
\item 6/11
\item 16/11
\end{enumerate}

\item $C$ sig‘imli ikki kondensator $2U$ va $U$ kuchlanish bilan zaryadlangan. Agar kondensator qutblari teskari holda ulansa, natijaviy kuchlanish qanday bo‘ladi?
\begin{enumerate}[label=\Alph*)]
\item 4U
\item 2,5U
\item 2U
\item 0,5U
\end{enumerate}

\item $10^4 \, nF$ sig‘imli ikki kondensator 5 V va 10 V kuchlanish bilan zaryadlangan. Ularning qutblari mos ravishda ulansa, natijaviy kuchlanish qanday bo‘ladi?
\begin{enumerate}[label=\Alph*)]
\item 5
\item 12,5
\item 7,5
\item 50
\end{enumerate}

\item Bir xil zaryadlangan $C$ va $2C$ sig‘imli kondensatorlar o‘zaro parallel ulansa, birinchi kondensatordagi kuchlanish qanday o‘zgaradi?
\begin{enumerate}[label=\Alph*)]
\item o‘zgarmaydi
\item 1,5 marta ortadi
\item 1,5 marta kamayadi
\item 3 marta ortadi
\end{enumerate}

\item 800 V potensiallar farqigacha zaryadlangan havo kondensatori o‘lchamlari birday bo‘lgan dielektrikli zaryadlanmagan kondensator bilan parallel ulangan. Potensiallar farqi 100 V bo‘lsa, dielektrikning singdiruvchanligini toping.
\begin{enumerate}[label=\Alph*)]
\item 7
\item 2
\item 16
\item 3
\end{enumerate}

\item Yassi havo kondensatori 50 V gacha zaryadlangan va manbadan uzilgan. Qoplamalar orasiga 1 mm qalinlikdagi metall list kiritildi. Masofa 5 mm bo‘lsa, yangi potensiallar farqini aniqlang (V).
\begin{enumerate}[label=\Alph*)]
\item 30
\item 15
\item 80
\item 40
\end{enumerate}

\item Ikkita bir xil kondensator parallel ulangan va 6 V gacha zaryadlangan. Uzilgach, birining plastinkalari orasidagi masofa 2 marta kamaytirilsa, yangi potensiallar farqi qancha (V)?
\begin{enumerate}[label=\Alph*)]
\item 6
\item 2
\item 4
\item 12
\end{enumerate}

\item 12 $\mu F$ sig‘imli kondensatorga parallel voltmetr 3 V ni ko‘rsatmoqda. Unga 8 $\mu F$ sig‘imli ikkinchi kondensator parallel ulansa, voltmetr necha V ni ko‘rsatadi?
\begin{enumerate}[label=\Alph*)]
\item 1,5
\item 3
\item 1
\item 2
\end{enumerate}

\end{enumerate}



\begin{table}[h!]
\centering
\begin{tabular}{|c|c|c|c|c|c|c|c|c|c|}
\hline
№ & J. & № & J. & № & J. & № & J. & № & J. \\ \hline
1 & D & 6 & A & 11 & A & 16 & D & 21 & C \\ \hline
2 & B & 7 & D & 12 & A & 17 & A & 22 & A \\ \hline
3 & D & 8 & A & 13 & D & 18 & B & 23 & D \\ \hline
4 & A & 9 & B & 14 & D & 19 & D & 24 & C \\ \hline
5 & D & 10 & D & 15 & D & 20 & C & 25 & B \\ \hline
\end{tabular}
\end{table}